\section{The Skellam rule for all loads}\label{sec:inversion}
The Gaussian rule's conditional miss rate guarantee requires $d\gg\sqrt M$. If we size the IBLT table for the expected load, $d$ can be far below $\sqrt M$. Failure is then rare, but the two-round protocol must still use the post-failure observation to size the second-round table. In this range, the conditional miss rate of the Gaussian rule can be several times the target $u$ (Table~\ref{tab:gauss-c} in the appendix).

The reason lies in the shape of the post-failure error at low load. Failure then has essentially one source: two differing keys select exactly the same $k$ cells. These cells contain only these two keys. The table has no pure cell, so peeling cannot start. Under random signs, the two keys are equally likely to have the same or opposite signs. If we place each key in the table alone, it leaves a count of $\pm1$ in each of its $k$ cells. Together, their squared counts sum to $2k$. Dividing by $k$ gives their exact contribution of $2$. With opposite signs, the counts in all $k$ shared cells cancel to $0$. The entire sum of squares of $2k$ disappears, leaving the estimate $2$ below the true value. With equal signs, each shared cell has count $\pm2$ and squared count $4$. The sum of squares becomes $4k$, leaving the estimate $2$ above the true value. This jump is equally likely to be upward or downward, and its size is independent of $M$. The remaining keys also occasionally share a cell, almost always exactly one. Each such overlap shifts the estimate by $2/k$, again with either direction equally likely. The number of such key pairs is approximately Poisson\cite{arratia1989chen,barbour1992poisson} and remains finite when $c$ is bounded. The post-failure error thus consists of one large jump of size $2$ and several small jumps of size $2/k$.

The Gaussian approximation cannot capture these two types of jumps. When $d$ is small relative to $\sqrt M$, the standard deviation $s(d)$ of the Gaussian rule tends to zero. The downward jump of size $2$ remains unchanged and lies arbitrarily many standard deviations into the Gaussian tail. The probability of rejecting the true value therefore tends to one half. When $c$ is of order one, only finitely many small jumps occur. The error still has a lattice distribution, with a visible departure from the Gaussian in the lower tail corresponding to the target risk (Appendix~\ref{app:limit}). We construct a null distribution that retains both types of jumps and approaches the Gaussian used by the Gaussian rule as $c$ grows (Figure~\ref{fig:aux}). The resulting Skellam rule controls the conditional miss rate for all $d\ge2$. Appendices~\ref{app:inversion} and~\ref{app:search} give the full proof and numerical implementation.

\subsection{The lattice-preserving test statistic \texorpdfstring{$\widetilde d$}{d~}}\label{sec:raw}
To keep these jumps exactly on lattice points, we need a discrete test statistic. The correction term $(\sum_vC_v)^2/M$ in $\widehat d$ shifts its values off the lattice. We therefore use the normalized sum of squares without subtracting the mean term,
\begin{equation}
\widetilde d=\frac1k\sum_{v=1}^M C_v^2.\label{eq:raw-energy}
\end{equation}
Expanding the sum of squares gives $d$ from the diagonal terms, leaving only overlap contributions to the error:
\begin{equation}
\widetilde d-d=\frac2k\sum_{i<j}H_{ij}\,\varepsilon_i\varepsilon_j.\label{eq:raw-expand}
\end{equation}
Every term on the right is an integer multiple of $2/k$, so $\widetilde d-d$ takes only integer multiples of $2/k$. The expansion shows both types of jumps directly. A duplicate edge pair has $H_{ij}=k$ and contributes exactly $2\varepsilon_i\varepsilon_j$, or $\pm2$ with equal probability. An edge pair sharing a single cell has $H_{ij}=1$ and contributes $\pm2/k$.

Without the correction term, $\widetilde d$ remains exactly unbiased for every mapping.

\begin{proposition}[Conditional unbiasedness of $\widetilde d$]\label{prop:raw-unbiased}
For every fixed mapping $\mathcal G$, $\E[\widetilde d\mid\mathcal G]=d$. Thus, for any mapping event $\mathcal S$ with positive probability, $\E[\widetilde d\mid\mathcal S]=d$. In particular, $\E[\widetilde d\mid F]=d$ holds for all finite $M$ and all sign compositions, however rare the failure.
\end{proposition}

\begin{proof}
Given $\mathcal G$, $H_{ij}$ is constant. By Proposition~\ref{prop:composition}, the signs remain iid and uniform, so $\E[\varepsilon_i\varepsilon_j\mid\mathcal G]=0$ for $i\ne j$. Taking expectations term by term in~\eqref{eq:raw-expand} gives the result.
\end{proof}

This property also follows from the random signs. In the standard IBLT, the cross terms have fixed sign products $s_is_j$. The same expansion gives the unconditional mean of the sum of squares as $d+\frac kM(b^2-d)$, which depends on the unknown $b$. The correction term in $\widehat d$ removes this dependence. With random signs, the cross terms have mean zero for every mapping. The sum of squares is therefore unbiased without the correction term.

$\widetilde d$ and $\widehat d$ differ only by the squared sum of the counts, weighted by $k/M$. This correction depends only on the random signs and is the classical AMS second-moment estimator\cite{alon1999spacecomplexity} (Appendix~\ref{app:ams}). After standardization by $s(d)$, the difference tends to zero, so Theorem~\ref{thm:clt} also holds for $\widetilde d$ (Appendix~\ref{app:clt}).

\subsection{Constructing the null distribution from the two jumps}\label{sec:aux}
For each candidate $n$, we construct a reference distribution for $\widetilde d$ after failure when the true value is $n$. This is the null distribution, which we use in Section~\ref{sec:skellam} to test each candidate. We construct it by counting the occurrences of the two types of jumps.

The large jump occurs exactly once and contributes $2\varepsilon_1\varepsilon_2$. Its direction is determined by the sign product of the duplicate edge pair.

The small jumps come from the $n-2$ keys outside the duplicate edge pair. They form $\binom{n-2}2$ pairs. Each pair shares a cell with probability approximately $k^2/M$, producing a small jump. Upward and downward jumps are equally likely, so each has expected count
\begin{equation}
\lambda_n=\frac12\binom{n-2}2\frac{k^2}M=\frac{k^2(n-2)(n-3)}{4M}.\label{eq:lambda}
\end{equation}
Key pairs sharing a cell are sparse. By Poisson splitting\cite{kingman1993poisson}, the upward and downward counts are approximately independent $\Pois(\lambda_n)$ counts. The upward count minus the downward count gives the net number of steps and follows a Skellam distribution\cite{skellam1946}, denoted $Z_n\sim\Skellam(\lambda_n,\lambda_n)$. Each step has size $2/k$.

Adding both types of jumps to the candidate $n$ gives the null distribution.
\begin{definition}[null distribution]\label{def:null}
For $\lambda\ge0$, let
\begin{equation}
J_\lambda=2\varepsilon_1\varepsilon_2+\frac2kZ,\qquad Z\sim\Skellam(\lambda,\lambda),\label{eq:J}
\end{equation}
where $\varepsilon_1\varepsilon_2$ takes the values $\pm1$ with equal probability and is independent of $Z$. We call $2\varepsilon_1\varepsilon_2$ the \textit{$\pm2$ term} and $\frac2kZ$ the \textit{Skellam term}. For a candidate integer $n\ge2$, we take $\lambda_n$ from~\eqref{eq:lambda}. The null distribution for candidate $n$ is the distribution of
\begin{equation}
n+J_{\lambda_n}=n+2\varepsilon_1\varepsilon_2+\frac2kZ_n\label{eq:aux}
\end{equation}
with CDF denoted by $G_n$.
\end{definition}
In Section~\ref{sec:uniform}, we prove that the post-failure distribution of $\widetilde d$ approaches $G_d$ uniformly over all $d\ge2$ (Theorems~\ref{thm:cross} and~\ref{thm:uniform}).

When $n=2,3$, fewer than two keys remain, so $\lambda_n=0$ and the null distribution assigns probability one half to each of $n\pm2$. In general, its mean is $n$. Since $\Var Z_n=2\lambda_n$, its variance is $4+2(n-2)(n-3)/M$.

\begin{figure}[!htbp]
\centering\includegraphics[width=\linewidth]{../figure/fig-aux.pdf}
\caption{Null distributions at three values of $c=d/\sqrt M$ ($k=3$, with $\lambda_d$ set to its limit $k^2c^2/4$). Blue and orange correspond to opposite and equal signs of the duplicate edge pair. The dashed line is the Gaussian used by the Gaussian rule.}
\label{fig:aux}
\end{figure}

The shape of this family explains why it applies at all loads. Figure~\ref{fig:aux} shows how the null distribution varies with $c$. When $c$ is small, the Skellam term is almost zero and the distribution is concentrated at $\pm2$. The Gaussian used by the Gaussian rule concentrates between these points, which lie beyond its tails. When $c\approx1$, the opposite-sign and equal-sign cases each spread into a lattice distribution with spacing $2/k$. These distributions overlap, but their combined shape still differs visibly from the Gaussian. As $c$ grows further, the width of the Skellam term far exceeds $2$, and the $\pm2$ term becomes a small contribution. The null distribution then approaches the Gaussian used by the Gaussian rule. This family thus connects the small-$c$ region to the region where the Gaussian rule applies. Section~\ref{sec:uniform} gives the rigorous approximation guarantee.

\subsection{The null distribution approximates the post-failure distribution at all loads}\label{sec:uniform}
Definition~\ref{def:null} constructs the null distribution from a counting argument. To use it to test candidates, we establish that it describes the post-failure distribution of $\widetilde d$ for all $d$. The receiver does not know the load region containing $d$ and cannot use the load to select a rule. The same family $G_d$ must therefore apply for both small and large $d$. We analyze these two cases and then combine them.

When $d$ is of the same order as $\sqrt M$ or smaller, overlaps are sparse. The model underlying $G_d$ consists of a duplicate edge pair and scattered single-cell overlaps among the remaining keys. The actual mapping may have additional structure: the duplicate edge pair may share a cell with a remaining key, two keys may share more than one cell, or one key may share cells with two other keys. Each case is much rarer than a single-cell overlap, and its probability tends to zero as $M$ grows. After excluding these cases, the actual mapping has exactly the structure underlying $G_d$. The actual distribution is therefore close to $G_d$ at every lattice point (Theorem~\ref{thm:cross}), as illustrated by the first two shapes in Figure~\ref{fig:aux}.

When $d\gg\sqrt M$, intersecting edge pairs are numerous and interwoven, so the sparse structure above no longer applies. The null distribution $G_d$ still applies because both distributions approach the same Gaussian. The statistic $\widetilde d$ is a sum of many small terms, each with equally likely signs, and approaches a Gaussian by Theorem~\ref{thm:clt}. The null distribution $G_d$ also approaches a Gaussian, as the third shape in Figure~\ref{fig:aux} shows, with width matching $s(d)$. The two distributions differ in detail but share the same Gaussian limit. Theorem~\ref{thm:uniform} combines these cases to give a uniform approximation over all $d\ge2$.

\begin{theorem}[Discrete limit distribution when $c$ is bounded]\label{thm:cross}
Fix $k\ge3$ and a finite constant $C$.
\begin{enumerate}[label=(\roman*)]
\item For $2\le d\le C\sqrt M$, the total variation distance between the conditional distribution of $\widetilde d$ given $F$ and the null distribution $G_d$ tends to zero uniformly and is at most
\begin{equation}
O_k\Bigl(\frac dM+\frac{d^3}{M^2}\Bigr).\label{eq:tv-error}
\end{equation}
\item If $d/\sqrt M\to c\in(0,\infty)$, the conditional distribution of $\widetilde d-d$ given $F$ converges in total variation to the distribution of $J_{k^2c^2/4}$, whose variance is $4+2c^2$.
\item The limiting distributions in (ii) satisfy
\begin{equation*}
J_{k^2c^2/4}\Rightarrow2\varepsilon_1\varepsilon_2\quad(c\to0),\qquad\frac{J_{k^2c^2/4}}{\sqrt2\,c}\Rightarrow\mathcal N(0,1)\quad(c\to\infty).
\end{equation*}
\end{enumerate}
\end{theorem}

Figure~\ref{fig:crosscdf} compares the lower tail of this discrete limit distribution with the Gaussian approximation, using planted samples with $c$ ranging from 0.25 to 8. Section~\ref{sec:exp-dist} also reports the distance between the empirical and limit distributions at fixed $c$ as $M$ grows.

\begin{proof}[Proof sketch]
We prove (i) in two steps. First, we replace the actual failure with a duplicate edge pair. By Proposition~\ref{prop:plant}, we can replace the mapping conditional on failure with the planted distribution, incurring error $O_k(d/M^2+M^{1-k})$ in the range of the theorem (Lemma~\ref{lem:plant}). Second, we exclude the three cases above from the planted mapping. Their probabilities are $O_k(d/M)$, $O_k(d^2/M^2)$, and $O_k(d^3/M^2)$, respectively. The remaining structure consists of a duplicate edge pair and mutually disjoint single-cell overlaps. The duplicate edge pair gives the $\pm2$ term. By the Chen--Stein local dependence bound\cite[Theorem~1]{arratia1989chen}, the overlap count is approximately $\Pois(2\lambda_d)$. Independent signs split this count into two independent $\Pois(\lambda_d)$ counts. Their difference gives the Skellam term, so $\widetilde d-d$ is close to $J_{\lambda_d}$. The Poisson approximation has error $O_k((d^2+d^3)/M^2)$. Since $2\le d\le M$ and $k\ge3$, each of $d/M^2$, $d^2/M^2$, and $M^{1-k}$ is at most $d/M$. Also, $(d^2+d^3)/M^2$ is at most $2d^3/M^2$. Combining these bounds gives~\eqref{eq:tv-error}.
(ii) By (i), it suffices to show that $J_{\lambda_d}$ converges in total variation to $J_{k^2c^2/4}$. When $d/\sqrt M\to c$, we have $\lambda_d\to k^2c^2/4$. The required convergence follows from continuity in the Poisson parameter. Independence of the two terms gives $\Var J_\lambda=4+8\lambda/k^2$. Substituting $\lambda=k^2c^2/4$ gives $4+2c^2$.
(iii) As $c\to0$, $\lambda\to0$ and the Skellam term tends to zero in probability, leaving only the $\pm2$ term of $J_\lambda$. As $c\to\infty$, $\Var J_{k^2c^2/4}\sim2c^2$ and the Skellam term satisfies the CLT. The $\pm2$ term divided by $\sqrt2\,c$ tends to zero. Appendix~\ref{app:inversion} gives the full argument.
\end{proof}

\begin{theorem}[Uniform CDF approximation]\label{thm:uniform}
Fix $k\ge3$. The Kolmogorov distance between the conditional distribution of $\widetilde d$ given $F$ and the null distribution $G_d$ satisfies
\begin{equation}
\sup_{d\ge2}\ \sup_{x\in\mathbb R}\bigl|\Pr(\widetilde d\le x\mid F)-G_d(x)\bigr|\longrightarrow0
\quad(M\to\infty).\label{eq:uniform}
\end{equation}
\end{theorem}

\begin{proof}[Proof sketch]
When $d\gg\sqrt M$, Gaussian limits let us handle a dense intersection graph without extending the matching structure. By Theorem~\ref{thm:clt}, the conditional distribution of $\widetilde d$, standardized by $s(d)$, approaches the standard Gaussian. The null distribution is that of $d+J_{\lambda_d}$. When $\lambda_d\to\infty$, the Skellam term of $J_{\lambda_d}$ satisfies the CLT. The ratio of $\Var J_{\lambda_d}=4+2(d-2)(d-3)/M$ to $s(d)^2$ tends to one, and the $\pm2$ term divided by $s(d)$ tends to zero. Both distributions approach the same continuous Gaussian CDF, so their Kolmogorov distance tends to zero. This argument applies to $d=o(M)$, a positive load ratio, and unbounded overload. We obtain uniformity by a subsequence argument. Take a subsequence along which $d/\sqrt M$ converges. Theorem~\ref{thm:cross} handles finite limits, and the Gaussian argument above handles infinite limits (Appendix~\ref{app:inversion}).
\end{proof}

The same family of null distributions $G_d$ therefore describes the post-failure distribution of $\widetilde d$ for all $d\ge2$. For small $d$, the approximation holds pointwise on the lattice. For large $d$, both distributions approach the same Gaussian. We can therefore use this family without knowing the load region. In Section~\ref{sec:skellam}, we use it to construct the Skellam rule.

\subsection{The Skellam rule controls the conditional miss rate at all loads}\label{sec:skellam}
The Skellam rule uses the test inversion of the Gaussian rule (Section~\ref{sec:gauss}), replacing the Gaussian lower tail with the lattice lower tail of $G_n$. If the observation falls in the lower $u$ tail of $G_n$, candidate $n$ is too large to be the true value, so we reject candidate $n$. We take $U$ to be the largest accepted candidate.

\begin{definition}[Skellam rule]\label{def:skellam}
Let $\widetilde d$ be the observation and $u\in(0,1/2)$ the target risk. For each candidate integer $n\ge2$, we accept $n$ if $G_n(\widetilde d)>u$ and reject it otherwise. We define the upper bound as the supremum of all accepted candidates,
\begin{equation}
U=\sup\bigl\{n\ge2:\ G_n(\widetilde d)>u\bigr\},\label{eq:skellam-U}
\end{equation}
with $U=+\infty$ when the accepted candidates are unbounded.
\end{definition}

The null distribution for candidate $2$ takes only the values $0$ and $4$. Since $\widetilde d\ge0$, we always accept candidate $2$, so $U\ge2$. For practical table sizes, $U$ is also finite (for $u=0.01$, $M>18$ suffices). We can compute it exactly by a finite search, at a cost roughly proportional to $\lambda_d$ per computation (Appendix~\ref{app:search}).

The behavior at small and large loads explains why the Skellam rule applies throughout. When $c$ is small, the Skellam term is almost zero. The null distribution $G_n$ assigns half its mass to $n-2$ and half to $n+2$. We accept candidate $n$ exactly when the smaller value $n-2$ does not exceed the observation, giving $U\approx\widetilde d+2$. The rule thus adds back a potentially lost large jump. When the duplicate edge pair has opposite signs, the observation is $d-2$ and $U$ equals $d$ exactly. With equal signs, the observation is $d+2$ and $U$ is about $d+4$. Neither case produces a miss. In this region, the scale $s(n)$ of the Gaussian rule tends to zero, so its upper bound $U$ stays close to the observation. It underestimates $d$ when the duplicate edge pair has opposite signs, giving a conditional miss rate that tends to one half. When $c$ is large, the Skellam term dominates and $G_n$ approaches a Gaussian. The Skellam rule then approaches the Gaussian rule.

Theorem~\ref{thm:coverage} gives the main result of this section. For sufficiently large tables, whatever load region contains $d$, the probability that the Skellam rule underestimates $d$ after failure is at most the target risk $u$. This guarantee holds for any sign composition. The rule uses only $\sum_vC_v^2$ from the same table and the known parameters $(M,k,u)$. The Gaussian rule provides this guarantee only when $d\gg\sqrt M$ (Corollary~\ref{cor:gauss}). In Section~\ref{sec:protocol}, we apply this guarantee to the risk decomposition of the two-round protocol.

\begin{theorem}[Uniform asymptotic upper bound on the conditional miss rate]\label{thm:coverage}
Fix $k\ge3$ and $u\in(0,1/2)$. For any sign composition, the Skellam rule satisfies
\begin{equation}
\limsup_{M\to\infty}\ \sup_{d\ge2}\ \Pr(U<d\mid F)\le u.\label{eq:coverage}
\end{equation}
\end{theorem}

\begin{proof}[Proof sketch]
If the observation has distribution $G_n$ exactly, we reject $n$ when it falls in a lower tail with total probability at most $u$. By the definition of the CDF, the rejection probability is at most $u$, including for discrete distributions. The rejection region is a lower set, so the rejection probabilities under the actual distribution and $G_n$ differ by at most their Kolmogorov distance. A miss implies rejection of the true value: $\{U<d\}\subseteq\{G_d(\widetilde d)\le u\}$. Taking $n=d$ and applying Theorem~\ref{thm:uniform} gives~\eqref{eq:coverage}. By Proposition~\ref{prop:composition}, the joint distribution of the coding coefficients is independent of the sign composition. The result therefore holds for any sign composition. Appendix~\ref{app:inversion} gives the full proof.
\end{proof}

Theorem~\ref{thm:coverage} is asymptotic in the table size. Two effects arise on finite tables. First, $G_n$ is a lattice distribution, so the rejection region contains only whole probability masses. The actual conditional miss rate can therefore be strictly below $u$. Second, the null distribution omits overlaps between the duplicate edge pair and the remaining keys. These occur at a rate of about $(d-2)k^2/M$ and can push the conditional miss rate slightly above $u$ on small tables. Section~\ref{sec:exp-cover} evaluates finite-table behavior experimentally (Figures~\ref{fig:coverage} and~\ref{fig:confirm}). At $M=64$, the largest miss/u is 1.08, an excess beyond sampling error. All configurations with $M\ge256$ meet the target within sampling error. The effect of the omitted overlaps decreases as $1/M$, and extra capacity also decreases as $M$ grows. When the first-round table must be small, the second-round sizing margin $1/(0.8\alpha_c)$ absorbs slight underestimation (Section~\ref{sec:protocol}). If more margin is needed, we can reduce the target load from 0.8, using a larger second-round table to lower the two-round failure rate.

\subsection{Combining the two rules using the first-round observation}\label{sec:switch}
The two rules have limitations in different ranges. The Gaussian rule underestimates $d$ only when $c$ is bounded, where the Skellam rule is inexpensive to compute. The cost of the Skellam rule grows roughly with $\lambda_d$, or $c^2$, while the Gaussian rule meets the target for large $c$. A system can therefore combine the rules using the first-round observation. We use the Gaussian rule when $\widetilde d$ is at or above a switch point and the Skellam rule otherwise. We compute $\widetilde d$ directly from the first-round counts, so switching adds no cost.

The combined rule retains the uniform guarantee of Theorem~\ref{thm:coverage} if the switch point grows faster than $\sqrt{M_1}$ as $M_1$ increases (Appendix~\ref{app:switch}). When $d/\sqrt{M_1}$ is bounded, $\widetilde d$ falls below the switch point with probability tending to one, so we use the Skellam rule. When $d/\sqrt{M_1}\to\infty$, both rules compare an asymptotically Gaussian statistic with the same lower $u$ quantile and give asymptotically identical conclusions. For example, a switch point of $1.5\alpha_cM_1$ corresponds to a first-round load of 1.5. Above load 1.5, sizing requires only one monotone binary search, and the protocol takes the same time in the overload region as with the Gaussian rule (Section~\ref{sec:exp-cost}).

\subsection{From the conditional miss rate to the failure probability of the two-round protocol}\label{sec:protocol}
We now connect the post-failure conditional miss rate to the overall success rate of two-round reconciliation. The protocol follows Section~\ref{sec:flow}. We denote first-round failure by $F_1$ and second-round failure by $F_2$. Let $E$ denote a checksum error: a non-pure cell passes the check, a key that passes the check lies outside the set difference, or all fields are zero at termination despite a nonempty core.

\begin{theorem}[Two-round risk decomposition]\label{thm:protocol}
Under the ideal hash model, let $\Pr_m(F)$ denote the probability of peeling failure for $d$ differing keys in a fresh table with $m$ cells. If the reserve size is at least $m(d)$, then
\begin{equation}
\Pr\big((F_1\cap F_2)\cup E\big)
\le\Pr(F_1,\ U_b<d)+\Pr(F_1)\max_{m\ge m(d)}\Pr\nolimits_m(F)+\Pr(E).\label{eq:risk}
\end{equation}
\end{theorem}

\begin{proof}
By the union bound, it suffices to bound $F_1\cap F_2$ and $E$ separately. The upper bound $U_b$ is a function of the first-round counts and the set sizes. The second-round mapping is independent of it. Under ideal purity, $F_2$ depends only on the second-round mapping and $d$. Hence $\Pr(F_1\cap F_2)=\E[\mathbf 1_{F_1}\Pr_{m(U_b)}(F)]$. On $\{U_b<d\}$, we bound the conditional failure probability by one. On $\{U_b\ge d\}$, we have $m(U_b)\ge m(d)$. Taking the supremum gives~\eqref{eq:risk}.
\end{proof}

The three terms correspond to three sources of error. The first is a sizing error: after first-round failure, the upper bound underestimates $d$. The second-round table is sized using $U_b$ and has insufficient capacity. This term is the joint miss rate $\Pr(F_1)\Pr(U_b<d\mid F_1)$. For the Skellam rule, Theorem~\ref{thm:coverage} bounds it asymptotically by $u\Pr(F_1)$, uniformly over all $d\ge2$. For a one-sided difference, $U_b\ge|b|=d$, so this term is zero. The Gaussian rule can have a conditional miss rate far above $u$ when $c$ is bounded, although failure itself is then rare. By Proposition~\ref{prop:plant}, when $d=o(M_1)$, $\Pr(F_1)$ is asymptotically equal to the expected number of duplicate edge pairs, $\binom d2/\binom{M_1}k$. This is about $2.8\times10^{-6}$ for $k=3$, $M_1=1024$, and $d=\sqrt{M_1}=32$, so the joint miss rate remains small.

The second term comes from the randomness of peeling. Even with sufficient second-round capacity, the random hypergraph can leave a nonempty 2-core with small probability. The second-round load is at most $0.8\alpha_c$, so Lemma~\ref{lem:propL} gives a failure probability that decays polynomially in $d$. The third term is a verification error: the checksum identifies a non-pure cell as pure or indicates termination on a nonempty core, leading to incorrect recovery. Each check passes falsely with probability at most $2^{-64}$. The number of checks is at most the number of cells plus $k$ times the number of deleted keys. The keySum and count fields of a nonempty core can both cancel to zero. The probability that all checksums are then also zero is at most $2^{-64}$, and this termination error can occur at most once per round. Thus $\Pr(E)$ is of order $2(M_1+m+kd)\cdot2^{-64}$, which is negligible at practical scales. With 32-bit checksums, this term is of order $2(M_1+m+kd)\cdot2^{-32}$, at most $10^{-4}$ for $M_1=16384$ and load at most 3. Since $\widetilde d$ uses only the count field, checksum width does not affect the first term. At deployment, we allocate the target failure probability between the sizing and checksum terms.

The bound in~\eqref{eq:risk} adds the three terms and conservatively treats every miss in the first term as a two-round failure. The second-round table has spare capacity, so it still succeeds with high probability under slight underestimation. The actual contribution of such misses is on the same scale as the second term. The protocol experiments in Section~\ref{sec:exp-e2e} also show that two-round failures arise almost entirely from the second term. Both rules have two-round failure rates comparable to that of the oracle upper bound ($U=d$).
